% Exported from Math++. Compile with pdflatex.
\documentclass[11pt,a4paper]{article}
\usepackage[T1]{fontenc}
\usepackage[utf8]{inputenc}
\usepackage{lmodern}
\usepackage{amsmath,amssymb,amsthm,mathtools}
\usepackage{graphicx,booktabs,xcolor}
\usepackage[hidelinks]{hyperref}
\usepackage[margin=25mm]{geometry}
\providecommand{\degree}{\ensuremath{^\circ}}
\providecommand{\xmark}{\ensuremath{\times}}
\providecommand{\omicron}{o}
\providecommand{\oiint}{\mathop{\oint\mkern-12mu\oint}}
\theoremstyle{plain}
\newtheorem{theorem}{Theorem}
\newtheorem{lemma}[theorem]{Lemma}
\newtheorem{corollary}[theorem]{Corollary}
\newtheorem{proposition}[theorem]{Proposition}
\newtheorem{conjecture}[theorem]{Conjecture}
\theoremstyle{definition}
\newtheorem{definition}[theorem]{Definition}
\newtheorem{example}[theorem]{Example}
\theoremstyle{remark}
\newtheorem*{remark}{Remark}
\setlength{\parindent}{0pt}\setlength{\parskip}{0.6em}
\title{Units in a worksheet}
\date{}
\begin{document}
\maketitle
A climber of mass $m \equiv 70\,\mathrm{kg}$ stands on Earth, where $g \equiv 9.81\,\mathrm{m/s^2}$.

Her weight is

\begin{equation*}
  W \coloneqq  m g
\end{equation*}
\begin{equation*}
  W = 686.7\,\mathrm{N}
\end{equation*}
A car drives at $v \coloneqq  100\,\mathrm{km/h}$ for $t \coloneqq  12\,\mathrm{s}$.

\begin{equation*}
  v \operatorname{in} \mathrm{m/s} = \frac{250}{9}\,\mathrm{m/s}
\end{equation*}
\begin{equation*}
  v t \operatorname{in} \mathrm{m} = \frac{1000}{3}\,\mathrm{m}
\end{equation*}
The room is at $T \coloneqq  25\,{}^\circ\mathrm{C}$; in Fahrenheit that is $T \operatorname{in} {}^\circ\mathrm{F} = 77\,\mathrm{^\circ F}$.

Angles in degrees are fine inside trigonometric functions: $\sin(30^\circ) = \frac{1}{2}$.

A $2\,\mathrm{kW}$ heater left on for $3\,\mathrm{h}$ uses $2\,\mathrm{kW} \cdot 3\,\mathrm{h} \operatorname{in} \mathrm{kWh} = 6\,\mathrm{kWh}$.

\section*{Mistakes the checker catches}

Adding a length to a time: $3\,\mathrm{m} + t = ?$

An exponential of a length: $\exp(2\,\mathrm{m}) = ?$

A letter is a unit only straight after a number or inside \texttt{\textbackslash{}mathrm\{\}}, so $m$ above is still the climber's mass.

\end{document}
